\section{Questão 8}
Seja $X$ uma variável aleatória contínua com a PDF abaixo, onde $\alpha$ é uma constante.
%
\begin{equation*}
    f_X(x) = 
    \begin{cases}
        0,             & x < 0 \\
        \alpha\pi x^2, & 0 \leq x \leq 1 \\
        0,             & x > 1
    \end{cases}
\end{equation*}

\begin{itemize}
    \item [a)] O valor de $\alpha$.
    \item [b)] Determine a CDF para a variável randômica.
    \item [c)] Determine $P(1.0 < x < 2.0)$.
    \item [d)] Determine $P(0.6 < x < 0.9)$.
    \item [e)] Determine $P((x \leq 0.5) \cup (x > 0.8))$.
    \item [f)] Calcule $\mu_X$.
\end{itemize}

\respostas

\begin{itemize}
    \item [a)] Pelos axiomas da probabilidade,
    \begin{equation}
        \int_\mathbb{R} f_X(x) \ dx = 1 \iff 
        \int_0^1 \alpha\pi x^2 \ dx = 1 \iff 
        %\alpha\pi \int_0^1 x^2 \ dx = 1 \iff 
        \alpha\pi \frac{x^3}{3}\bigg|_0^1 = 1 \iff 
        \alpha = \frac{3}{\pi}.
    \end{equation}
    
    \item [b)] Integrando a PDF,
    \begin{equation}
        F_X(x) = 
        \begin{cases}
            0,   & x < 0 \\
            x^3, & 0 \leq x \leq 1 \\
            1,   & x > 1
        \end{cases}
    \end{equation}

    \item [c), d), e)]
    Obtida a CDF do item (b), encontrar as probabilidades é direto.
    Denotando os intervalos de interesse por 
    $I_c = (1, 2)$, 
    $I_d = (0.6, 0.9)$ e 
    $I_e = (-\infty, 0.5] \cup (0.8, +\infty)$, tem-se
    
    \begin{minipage}{.28\textwidth}
        \begin{align*}
            P(I_1) 
            &= F(2) - F(1) \\
            &= 1^3 - 1^3 \\
            &= 0
        \end{align*}
    \end{minipage}
    %
    \begin{minipage}{.32\textwidth}
        \begin{align*}
            P(I_2) 
            &= F(0.9) - F(0.6) \\
            &= 0.9^3 - 0.6^3 \\
            &= 0.513
        \end{align*}
    \end{minipage}
    %
    \begin{minipage}{.36\textwidth}
        \begin{align*}
            P(I_3) 
            &= F(0.5) + 1-F(0.8) \\
            &= 0.5^3 + 1 - 0.8^3 \\
            &= 0.613
        \end{align*}
    \end{minipage}
    
    \item [f)] Por definição,
    \begin{equation}
        \mu_X = \E{X} 
        = \int_\mathbb{R} x f_X(x) \ dx
        = 3\int_\mathbb{R} x^3 \ dx
        = 3 \frac{x^4}{4}\bigg|_0^1
        = \frac{3}{4}.
    \end{equation}
\end{itemize}